\section{Parent Hamiltonians on finite PEPS graphs}%
\label{sec:peps_region_parent_hamiltonian}

The construction of \cite[Section~IV.C.1]{Cirac2021Matrix}, local source
lines 2003--2011, applies to arbitrary finite graphs and site-dependent
tensors. Different interaction terms may act on different regions.

\begin{definition}[Regional parent interaction]%
    \label{def:peps_region_parent_interaction}
    \lean{TNLean.PEPS.IsRegionParentInteraction}
    \leanok
    \uses{def:peps_region_ground_space}
    A parent interaction on a region $R$ is a positive semidefinite
    operator $h_R$ on its physical space satisfying
    $\ker h_R=\mathcal G_R$. It need not be an orthogonal projector.
\end{definition}

\begin{theorem}[Canonical regional parent interaction]%
    \label{thm:peps_canonical_region_interaction}
    \lean{TNLean.PEPS.regionGroundSpaceES,
        TNLean.PEPS.canonicalRegionParentInteraction,
        TNLean.PEPS.canonicalRegionParentInteraction_posSemidef,
        TNLean.PEPS.ker_canonicalRegionParentInteraction,
        TNLean.PEPS.isRegionParentInteraction_canonical}
    \leanok
    \uses{def:peps_region_parent_interaction}
    The orthogonal projector onto $\mathcal G_R^\perp$ is a regional
    parent interaction.
\end{theorem}

\begin{proof}\leanok
    Orthogonal projections are positive semidefinite, and the kernel
    of the projector onto $\mathcal G_R^\perp$ is
    $\mathcal G_R^{\perp\perp}=\mathcal G_R$.
\end{proof}

\begin{definition}[Parent Hamiltonian of a finite family of regions]%
    \label{def:peps_region_parent_hamiltonian}
    \lean{TNLean.PEPS.regionConfigEquiv,
        TNLean.PEPS.regionLocalTerm,
        TNLean.PEPS.regionSliceMap,
        TNLean.PEPS.regionParentGroundSpace,
        TNLean.PEPS.regionParentHamiltonian}
    \leanok
    \uses{def:peps_region_parent_interaction}
    Let $(R_i)_{i\in I}$ be a finite family of regions, and let $h_i$
    be a parent interaction on $R_i$. Its extension to the full graph
    and the parent Hamiltonian are
    \begin{align}
        \widetilde h_i & =h_i\otimes\mathbb 1_{R_i^c},
        \label{eq:peps_parent_extension}               \\
        H              & =\sum_{i\in I}\widetilde h_i.
        \label{eq:peps_parent_sum}
    \end{align}
    The corresponding space of regional boundary conditions is
    \begin{align}
        \mathcal K
         & =\{\phi: \phi(\cdot,\tau)\in\mathcal G_{R_i}
        \text{ for every }i\in I
        \text{ and every complementary configuration }\tau\}.
        \label{eq:peps_parent_common_kernel}
    \end{align}
\end{definition}

\begin{theorem}[Positivity and the common kernel]%
    \label{thm:peps_region_parent_common_kernel}
    \lean{TNLean.PEPS.regionLocalTerm_posSemidef,
        TNLean.PEPS.regionLocalTerm_mulVec_assemble,
        TNLean.PEPS.regionLocalTerm_mulVec_eq_zero_iff,
        TNLean.PEPS.regionParentHamiltonian_posSemidef,
        TNLean.PEPS.regionParentHamiltonian_mulVec_eq_zero_iff,
        TNLean.PEPS.ker_regionParentHamiltonian}
    \leanok
    \uses{def:peps_region_parent_hamiltonian}
    The parent Hamiltonian is positive semidefinite and satisfies
    \begin{align}
        \ker H & =\bigcap_{i\in I}\ker\widetilde h_i=\mathcal K.
        \label{eq:peps_parent_kernel}
    \end{align}
    Thus its kernel depends only on the regional PEPS spaces, rather
    than on the choice of positive interactions with those kernels.
\end{theorem}

\begin{proof}\leanok
    Positivity is preserved by tensoring with the identity and by
    finite sums. For each fixed complementary configuration $\tau$,
    (\ref{eq:peps_parent_extension}) acts on the slice
    $\phi(\cdot,\tau)$ as $h_i$. Finally,
    $\langle\phi,H\phi\rangle=\sum_i
        \langle\phi,\widetilde h_i\phi\rangle$ is a sum of nonnegative
    numbers. If $H\phi=0$, each summand is zero, and positivity gives
    $\widetilde h_i\phi=0$. The converse follows by summing.
\end{proof}

\begin{theorem}[Frustration freeness of the represented PEPS]%
    \label{thm:peps_region_parent_frustration_free}
    \lean{TNLean.PEPS.regionLocalTerm_stateCoeff_eq_zero,
        TNLean.PEPS.stateCoeff_mem_regionParentGroundSpace,
        TNLean.PEPS.regionParentHamiltonian_stateCoeff_eq_zero}
    \leanok
    \uses{def:peps_region_parent_hamiltonian,
        thm:peps_region_ground_space_slice}
    The closed PEPS vector $\psi$ satisfies
    $\widetilde h_i\psi=0$ for every $i\in I$, and hence $H\psi=0$.
    Whenever $\psi\ne0$, it is a ground state of this positive
    Hamiltonian, and it minimizes every local term simultaneously.
\end{theorem}

\begin{proof}\leanok
    \uses{thm:peps_region_ground_space_slice,
        thm:peps_region_parent_common_kernel}
    Every complementary slice of $\psi$ belongs to the regional
    PEPS ground space. Therefore every extended interaction kills
    $\psi$, and their sum does as well.
\end{proof}
